% TOP_PROBLEM: {"id": 30006774, "title": "Bochner-Riesz conjecture", "category_id": 8, "category": {"id": 8, "name": "analysis", "display_name": "Analysis", "description": "Limits, continuity, calculus, and function theory.", "slug": "analysis", "order_index": 8, "created_at": "2026-07-31T15:26:25.671Z"}, "status": "open"}
% FIRST_POSTED: 2026-09-27
% !TeX program = lualatex
% Compile twice: lualatex 30006774.tex
% All references and prose are embedded; no BibTeX database or external inputs.
\documentclass[11pt,a4paper]{article}
\usepackage[margin=24mm,headheight=14pt]{geometry}
\usepackage{fontspec}
\setmainfont{Latin Modern Roman}
\setsansfont{Latin Modern Sans}
\setmonofont{Latin Modern Mono}
\usepackage{amsmath,amssymb,mathtools}
\usepackage{unicode-math}
\setmathfont{Latin Modern Math}
\usepackage{microtype}
\usepackage{enumitem,xurl,needspace}
\usepackage[dvipsnames]{xcolor}
\usepackage{hyperref}
\hypersetup{unicode=true,colorlinks=true,linkcolor=MidnightBlue,urlcolor=MidnightBlue,
 pdftitle={500 Mathematical Problem Statements},pdfauthor={Source-annotated compilation},bookmarksnumbered=true}
\usepackage{bookmark}
\usepackage{tocloft}
\setlength{\cftsecnumwidth}{3em}
\cftsetpnumwidth{2.5em}
\cftsetrmarg{4em}
\usepackage{titlesec}
\titleformat{\section}{\Large\bfseries\raggedright}{\thesection}{0.7em}{}
\usepackage{fancyhdr}
\pagestyle{fancy}\fancyhf{}
\fancyhead[L]{\small\sffamily Mathematical problem statements}
\fancyhead[R]{\small Source edition 22}
\fancyfoot[C]{\thepage}
\setlength{\parindent}{0pt}
\setlength{\parskip}{5pt plus 1pt}
\setlength{\emergencystretch}{3em}
\setcounter{tocdepth}{1}
\setcounter{secnumdepth}{1}
\setlist[itemize]{leftmargin=3em,itemsep=5pt,topsep=4pt,parsep=0pt}
\allowdisplaybreaks
\newcommand{\N}{\mathbb{N}}
\newcommand{\Z}{\mathbb{Z}}
\newcommand{\Q}{\mathbb{Q}}
\newcommand{\R}{\mathbb{R}}
\newcommand{\C}{\mathbb{C}}
\newcommand{\F}{\mathbb{F}}
\newcommand{\A}{\mathbb{A}}
\newcommand{\PP}{\mathbb P}
\newcommand{\E}{\mathbb{E}}
\newcommand{\eps}{\varepsilon}
\newcommand{\dd}{\,\mathrm d}
\newcommand{\sref}[2]{\hyperlink{source:#1:#2}{[S#2]}}
\newcommand{\eref}[2]{\hyperlink{extra:#1:#2}{[E#2]}}
% Turn the next switch off for a shorter reading copy. The quotations preserve
% every exactTarget field, including qualifications omitted from a short summary.
\newif\ifshowcatalogue
\showcataloguetrue
\newcommand{\cataloguescope}[1]{\ifshowcatalogue
 \paragraph{Exact scope in the supplied catalogue.}
 \begin{quote}\small #1\end{quote}\fi}

% Standalone notebook wrapper; original problem section retained verbatim.
\numberwithin{equation}{section}
\hypersetup{pdftitle={Research attempt -- problem 30006774},
 pdfauthor={Research notebook; source compilation retained}}
\fancyhead[L]{\small\sffamily Research notebook}
\fancyhead[R]{\small\sffamily Problem 30006774 | 27 September 2026}
\begin{document}
\setcounter{section}{0}
% ================================================================
% problemId: 30006774
% Canonical problem ID: 30006774
% exactTarget SHA-256: 11ea0fb26f238812dd317e56615b3fa595cce1cd9349f2754c0f227fc4cc3a09
\clearpage
\section*{\texorpdfstring{Bochner-Riesz conjecture}{Bochner-Riesz conjecture}}\label{30006774}
\subsection{Definitions and mathematical statement}
For $R>0$, define the Euclidean Bochner--Riesz operator on Schwartz functions by
\[
 \widehat{S_R^\delta f}(\xi)=\left(1-\frac{|\xi|^2}{R^2}\right)_+^\delta\widehat f(\xi),
\]
where at $\delta=0$ the multiplier is the indicator of the ball, up to its irrelevant boundary. Uniform $L^p$ boundedness means
$\sup_{R>0}\|S_R^\delta\|_{L^p\to L^p}<\infty$.
For $d\ge2$ and $1<p<\infty$, the conjectured strong-type range is $\delta\ge0$ when $p=2$, and
$\delta>\max\{d|1/p-1/2|-1/2,0\}$ when $p\ne2$.
The problem is this full characterization in every dimension. Strong operator boundedness is distinct from weak-type endpoint estimates or almost-everywhere convergence, even when those topics appear in the source title.
\par\smallskip\textit{Formulation and concept references:} \sref{30006774}{1}.
\cataloguescope{Prove the sharp conjectured characterization of all orders delta greater than or equal to zero and exponents 1<p<infinity for which the Euclidean Bochner-Riesz means are uniformly bounded on L\textasciicircum{}p(R\textasciicircum{}d), in every dimension d at least 2.}
\subsection{Short English statement}
For exactly which summation orders and integrability exponents do spherical Fourier cutoffs remain uniformly bounded on Euclidean Lp spaces?
% BEGIN RESEARCH ADDITION ATTEMPT-180
\subsection{Research attempt}

\paragraph{Current frontier.}
\textbf{Editorial source-fit note.} The retained formulation reference
concerns almost-everywhere convergence on Heisenberg-type groups. That is
not the Euclidean strong-operator characterization requested here
\sref{30006774}{1}. For the latter, Guo--Oh--Wang--Wu--Zhang give an explicit
restriction-based sufficient range in Theorem 1.3 and explain a barrier to
treating the multiplier as a completely general positive-curvature
oscillatory integral \eref{30006774}{1}, Introduction. Their result supplies
substantial partial ranges, not all exponents in the question. Kim's
2024 manuscript, published in 2025, gives improved almost-everywhere
convergence conditions in dimensions two and three; that conclusion also
must not be substituted for uniform strong-type boundedness
\eref{30006774}{3}. The attempt below concentrates on the precise loss incurred
by estimating one thin spherical shell absolutely.

\paragraph{Attack route.}
A restriction transfer needs sharp estimates for the relevant oscillatory
operators rather than just their arbitrary-phase analogues. Polynomial
partitioning needs to control the narrow contributions without losing the
critical exponent. We pursue a simpler diagnostic: decompose at distance
$2^{-j}$ from the sphere, compute the exact loss of the elementary
kernel argument, and state the uniform shell improvement that would close
the positive part of the conjecture. The point is to isolate cancellation
that a proposed proof must actually deliver, not to hide the desired bound
inside an unspecified interpolation step.

\paragraph{Attempt: reduction to one radius and thin shells.}
If $D_Rf(x)=f(Rx)$, Fourier transformation gives
$S_R^\delta D_R=D_RS_1^\delta$. The $L^p$ scaling factors cancel on
conjugation, so the operator norm is independent of $R$. It suffices to
study $R=1$; this already handles the uniform-radius quantifier.

Fix a smooth dyadic partition on $0<t<1/4$. Up to a smooth multiplier
supported away from the unit sphere, write
\begin{equation}\label{eq180-shells}
 (1-|\xi|^2)_+^\delta
 =m_{\mathrm{int}}(\xi)+\sum_{j\ge j_0}2^{-j\delta}
          a_\delta\bigl(2^j(1-|\xi|^2)\bigr),
\end{equation}
where $a_\delta(s)=s^\delta\psi(s)$ with
$\psi\in C_c^\infty((1/2,2))$. Finitely many transition terms can be
absorbed in $m_{\mathrm{int}}$. For each fixed $\delta\ge0$, all the
rescaled profile derivatives are bounded independently of $j$. Constants
may depend on $\delta$; no uniformity as $\delta$ tends to an endpoint is
claimed. Let $P_j$ be the multiplier with the profile in the sum.

\textbf{Elementary shell bound.} For this fixed profile and
$t_p=|1/p-1/2|$,
\begin{equation}\label{eq180-elementary}
        \|P_j\|_{L^p\to L^p}\lesssim_{p,d,\delta}
                          2^{j(d-1)t_p}.
\end{equation}
Here is the kernel calculation identifying the loss. Put $h=2^{-j}$ and
$K_j=\mathcal F^{-1}(a_\delta(h^{-1}(1-|\xi|^2)))$. Polar coordinates
express it as an integral of the Fourier transform of surface measure
over a radial interval of length $O(h)$ around radius one. For $|x|\ge1$,
the spherical integral is a sum of the two phases
$e^{\pm2\pi i r|x|}|x|^{-(d-1)/2}$ with symbol amplitudes; their radial
derivatives obey the corresponding symbol bounds. Repeated integration
by parts in $r$ outside $h|x|\lesssim1$, and the absolute radial bound
inside, give, for every fixed $N$,
\[
 |K_j(x)|\le C_N h(1+|x|)^{-(d-1)/2}(1+h|x|)^{-N}.
\]
The bound for $|x|<1$ follows directly from the volume of the frequency
shell. Integrating first over $|x|\le h^{-1}$ and then its complement,
with $N>d$, yields
\[
        \|K_j\|_1\lesssim h^{-(d-1)/2}.
\]
Plancherel gives $\|P_j\|_{2\to2}\lesssim1$. Interpolation between
$L^1$ and $L^2$, followed by duality for $p>2$, proves
\eqref{eq180-elementary}. This uses the standard spherical
stationary-phase expansion underlying the reductions in
\eref{30006774}{1}, Section 2; it does not assume the restriction conjecture.

Summing the operator norms in \eqref{eq180-shells} proves the classical
sufficient condition
\begin{equation}\label{eq180-classical}
                     \delta>(d-1)|1/p-1/2|.
\end{equation}
The sum is genuinely convergent in operator norm, so no interchange of an
uncontrolled infinite series with the $L^p$ norm is involved.

\paragraph{Attempt: why the conjectured exponent is smaller.}
The full kernel is, up to a nonzero dimensional constant,
\[
 K^\delta(x)=|x|^{-d/2-\delta}
                         J_{d/2+\delta}(2\pi|x|).
\]
The Bessel asymptotic has leading term a nonzero multiple of
$|x|^{-(d+1)/2-\delta}\cos(2\pi|x|-\theta)$, with an error smaller by
one power of $|x|$. Choose a Schwartz function whose Fourier transform is
one on the closed unit ball. Its image is exactly $K^\delta$. Integrating
on periodic radial subintervals where the cosine is bounded away from
zero proves that this image is not in $L^p$ unless
\[
             \delta>\frac d p-\frac{d+1}{2}.
\]
At equality the radial integral diverges logarithmically; replacing the
oscillatory kernel by its average would incorrectly lose that obstruction.
Duality gives the corresponding condition at $p>2$. Combined, these yield
the familiar necessary strict bound when it is nonnegative,
\[
                 \delta>d|1/p-1/2|-\tfrac12.
\]
For the remaining $\delta=0$, $p\ne2$ cases, the ball-multiplier obstruction
is Fefferman's separate theorem \eref{30006774}{4}, not a consequence of this
integrability test. At $p=2$, all $\delta\ge0$ are bounded by Plancherel.

Set $\delta_c(p)=\max\{dt_p-1/2,0\}$. The loss in
\eqref{eq180-classical}, relative to the desired positive-order range, is
\begin{equation}\label{eq180-loss}
 (d-1)t_p-\delta_c(p)=
 \begin{cases}(d-1)t_p,&dt_p\le1/2,\\
               1/2-t_p,&dt_p\ge1/2.
 \end{cases}
\end{equation}
It is strictly positive for $1<p<\infty$, $p\ne2$. Thus neither optimizing
the dyadic summation nor keeping logarithmic factors more carefully can
repair the elementary argument: its error is a power of the shell width.

\paragraph{Attempt: an explicit sufficient missing estimate.}
For each fixed smooth profile of the kind in \eqref{eq180-shells}, consider
\begin{equation}\label{eq180-missing}
 \forall\eta>0\quad
 \|P_j\|_{p\to p}\le C_{p,d,\eta,a}\,
                         2^{j(\delta_c(p)+\eta)}
                  \quad(j\ge j_0).
\end{equation}
If this estimate holds for every $d,p$ under consideration, then every
$\delta>\delta_c(p)$ in the target follows. Indeed, fix $\delta$ first,
choose $0<\eta<\delta-\delta_c(p)$, apply the estimate to
$a=a_\delta$, and sum the convergent geometric series
$\sum_j2^{-j(\delta-\delta_c(p)-\eta)}$.
Together with the necessary results and the separate $p=2$ case this
would give the full stated characterization. This is a sufficient
reduction; no converse to this particular shell formulation is asserted.

The estimate is where angular interaction must enter. Applying absolute
values to the kernel before summation recovers only
\eqref{eq180-elementary}. A square-function argument that merely combines
already bounded shells does not, by itself, improve the estimate for a
single shell: it must also prove new information about the angular pieces
within that shell. Nor can one replace the phase by an arbitrary
positive-curvature phase and demand the same range. The sharpness examples
for that larger class are precisely a barrier explained in
\eref{30006774}{1}, Introduction.

\paragraph{Stress test.}
At $p=2$ the elementary shell exponent is already zero, but the sum of
norms at $\delta=0$ diverges. The actual ball projection is nevertheless
bounded by Plancherel. Thus failure of absolute summation is not evidence
of unboundedness; the endpoint must be handled separately. At $p\ne2$,
weak-type endpoint information and almost-everywhere convergence do not
supply the missing strong shell estimate. In particular the 2025 weighted
decoupling application \eref{30006774}{3} cannot be inserted into
\eqref{eq180-missing} simply by changing the words describing its norm.

There is also a dependency warning. The full Bochner--Riesz conjecture
implies the corresponding restriction conjecture by Tao's theorem
\eref{30006774}{2}. A proposed elementary cancellation argument proving
\eqref{eq180-missing} universally must therefore be checked against that
major additional consequence. No such cancellation has been proved here.

\paragraph{Outcome.}
\textbf{Outcome: Conditional reduction.}

The attempt verifies the uniform-radius reduction, derives the elementary
shell loss, and proves that \eqref{eq180-missing} would supply the entire
positive-order sufficient range. It also separates the two different
necessary obstructions and the $p=2$ endpoint. These are standard analytic
ingredients assembled and audited here, not a claimed new improvement to
the best Bochner--Riesz bounds.

What remains is the sharp shell estimate, uniformly in the shell scale.
The present calculation exposes exactly the power saving needed and why
absolute-kernel interpolation does not produce it. It neither proves the
full characterization nor constructs an example contradicting it.
% END RESEARCH ADDITION ATTEMPT-180
\subsection{Sources}
\begin{itemize}\raggedright
\item[\textnormal{[S1]}]\hypertarget{source:30006774:1}{} Almost everywhere convergence of Bochner-Riesz means on Heisenberg-type groups.
\url{https://londmathsoc.onlinelibrary.wiley.com/doi/full/10.1112/jlms.12401}.
{\small\textit{Catalogue formulation source.}}
% BEGIN RESEARCH ADDITION SOURCES-180
\item[\textnormal{[E1]}]\hypertarget{extra:30006774:1}{} Shaoming Guo, Changkeun
Oh, Hong Wang, Shukun Wu and Ruixiang Zhang, \emph{The Bochner--Riesz problem:
An old approach revisited}, arXiv:2104.11188, Introduction, Theorem 1.3,
and Section 2. \url{https://arxiv.org/html/2104.11188}.
{\small\textit{Consulted 27 September 2026; Euclidean operator bounds and
the arbitrary-phase barrier, not the group setting of [S1].}}
\item[\textnormal{[E2]}]\hypertarget{extra:30006774:2}{} Terence Tao,
\emph{The Bochner--Riesz conjecture implies the restriction conjecture},
Duke Mathematical Journal 96(2) (1999), 363--375.
\url{https://doi.org/10.1215/S0012-7094-99-09610-2}.
{\small\textit{Dependency theorem; also explicitly discussed in [E1],
Introduction. Bibliographic record checked 27 September 2026.}}
\item[\textnormal{[E3]}]\hypertarget{extra:30006774:3}{} Jongchon Kim,
\emph{Weighted decoupling estimates and the Bochner--Riesz means},
arXiv:2406.03741; Forum of Mathematics, Sigma 13 (2025), e167.
\url{https://arxiv.org/abs/2406.03741}.
{\small\textit{Consulted 27 September 2026; the cited application is
almost-everywhere convergence, not the full strong-type target.}}
\item[\textnormal{[E4]}]\hypertarget{extra:30006774:4}{} Charles Fefferman,
\emph{The multiplier problem for the ball}, Annals of Mathematics 94(2)
(1971), 330--336. \url{https://annals.math.princeton.edu/1971/94-2/p05}.
{\small\textit{Primary theorem at order zero; publication record checked
27 September 2026.}}
% END RESEARCH ADDITION SOURCES-180
\end{itemize}

\end{document}
